% Historical source (removed by author request): references/_OLD_Übungsaufgaben/Übungsblatt 3 - Vektorechnung/Blatt3.tex:18
\begin{ExerciseBlock}
\ExerciseHeading{01}{Betrag und Einheitsvektor}{vector_algebra_magnitude_and_unit_vector}
Gegeben seien die folgenden Vektoren
\begin{center}
$\vec{a}=\left(\begin{array}{c} 3 \\ 2 \\ -4 \end{array}\right)$, 
$\vec{b}=\left(\begin{array}{c} -2 \\ 0 \\ 4\end{array}\right)$ und
${\vec{c}}=\left(\begin{array}{c} -5 \\ 1 \\ 4 \end{array}\right)$. 
\end{center}

Berechnen Sie den Vektor $\vec{s}$, seinen Betrag $|\vec{s}|$ und den Einheitsvektor $\vec{e}_s$:
\begin{enumerate}[label=\alph*)]
	\item $\vec{s}=3\vec{a}-5\vec{b}+3\vec{c}$
	\item $\vec{s}=-2\left(\vec{b}+5\vec{c}\right)+5\left(\vec{a}-5\vec{b}\right)$
\end{enumerate} 

\end{ExerciseBlock}
\UseExerciseSolution{vector_algebra_magnitude_and_unit_vector}
